% GENERATED FILE. Edit canonical lessons, then run npm run generate:books.
% canonical-source-hash: fnv1a64:36bb366f0bd9d2c9
% canonical-lessons: TA-C365-katturai, TA-C365-talaippu, TA-C365-munnurai, TA-C365-curukkam, TA-C365-merkol

\chapter{Writing an Essay}
\label{ch:ta-writing-an-essay}
\hlchaptermodality{365}

\begin{chapteropening}
\textbf{By the end of this chapter:} \emph{I can talk about the parts of an essay.}

Complete the last lesson of chapter 365: I can talk about the parts of an essay.
\end{chapteropening}

\section[\texorpdfstring{\ta{கட்டுரை} (kaṭṭurai)}{கட்டுரை (kaṭṭurai)}]{\ta{கட்டுரை} (kaṭṭurai) — an essay; an article}

\label{lesson:TA-C365-katturai}



\begin{quote}
Before the new one: say the Tamil for to study, then the Tamil for admission.
\end{quote}

\subsection*{You'll want to know: \ta{கட்டுரை}}
\textbf{\ta{கட்டுரை}} — \emph{kaṭṭurai} — \textquotedblleft{}an essay; an article\textquotedblright{}.

\begin{quote}
\ta{என்} \ta{தங்கை} \ta{தன்} \ta{பள்ளியைப்} \ta{பற்றி} \ta{ஒரு} \ta{கட்டுரை} \ta{எழுதினாள்}.
\end{quote}

\begin{quote}
\emph{eṉ taṅkai taṉ paḷḷiyaip paṟṟi oru kaṭṭurai eḻutiṉāḷ.}
\end{quote}

\begin{quote}
My younger sister wrote an essay about her school.
\end{quote}

\ta{கட்டுரைப்} \ta{போட்டி}, an essay competition.

\begin{cousinweb}
\textbf{\ta{கட்டுரை}} (\emph{kaṭṭurai}) is native: \emph{kaṭṭu}, \textquotedblleft{}to tie, build\textquotedblright{} (DEDR 1147), and \emph{urai}, \textquotedblleft{}word, speech\textquotedblright{} (DEDR 648): words put together into a whole.
\end{cousinweb}

\begin{grammarlens}[title={What you've built}]
One more word for giving reasons and opinions.
\end{grammarlens}

\subsection*{Your turn}
\begin{itemize}
\raggedright
  \item \emph{Say it:} \emph{kaṭṭurai}
  \item \emph{Say it:} \emph{kaṭṭurai}, once more
  \item \emph{Say it:} \emph{eṉ taṅkai taṉ paḷḷiyaip paṟṟi oru kaṭṭurai eḻutiṉāḷ.} — My younger sister wrote an essay about her school.
  \item \emph{Recall:} say the Tamil for to study, then the Tamil for admission, then say \emph{kaṭṭurai} again
\end{itemize}

\begin{tcolorbox}[breakable,colback=teal!4,colframe=teal!35!black,title={Before you move on}]
What does \ta{கட்டுரை} mean? (\textquotedblleft{}An essay; an article\textquotedblright{}.) Say it once more.
\end{tcolorbox}
\section[\texorpdfstring{\ta{தலைப்பு} (talaippu)}{தலைப்பு (talaippu)}]{\ta{தலைப்பு} (talaippu) — a title; a heading}

\label{lesson:TA-C365-talaippu}



\begin{quote}
Before the new one: say the Tamil for admission, then the Tamil for an essay.
\end{quote}

\subsection*{You'll want to know: \ta{தலைப்பு}}
\textbf{\ta{தலைப்பு}} — \emph{talaippu} — \textquotedblleft{}a title; a heading\textquotedblright{}.

\begin{quote}
\ta{இந்தக்} \ta{கட்டுரையின்} \ta{தலைப்பு} \ta{என்ன}?
\end{quote}

\begin{quote}
\emph{intak kaṭṭuraiyiṉ talaippu eṉṉa?}
\end{quote}

\begin{quote}
What is the title of this essay?
\end{quote}

From \ta{தலை}, head: the head of a piece of writing.

\begin{cousinweb}
\textbf{\ta{தலைப்பு}} (\emph{talaippu}) is native Dravidian: DEDR 3103 lists Tamil \emph{talai}, \textquotedblleft{}head, top, end\textquotedblright{}, and \emph{talaippu}, \textquotedblleft{}beginning; edge of a cloth\textquotedblright{}, the noun that now names a heading.
\end{cousinweb}

\begin{grammarlens}[title={What you've built}]
One more word for giving reasons and opinions.
\end{grammarlens}

\subsection*{Your turn}
\begin{itemize}
\raggedright
  \item \emph{Say it:} \emph{talaippu}
  \item \emph{Say it:} \emph{talaippu}, once more
  \item \emph{Say it:} \emph{intak kaṭṭuraiyiṉ talaippu eṉṉa?} — What is the title of this essay?
  \item \emph{Recall:} say the Tamil for admission, then the Tamil for an essay, then say \emph{talaippu} again
\end{itemize}

\begin{tcolorbox}[breakable,colback=teal!4,colframe=teal!35!black,title={Before you move on}]
What does \ta{தலைப்பு} mean? (\textquotedblleft{}A title; a heading\textquotedblright{}.) Say it once more.
\end{tcolorbox}
\section[\texorpdfstring{\ta{முன்னுரை} (muṉṉurai)}{முன்னுரை (muṉṉurai)}]{\ta{முன்னுரை} (muṉṉurai) — a preface; an introduction}

\label{lesson:TA-C365-munnurai}



\begin{quote}
Before the new one: say the Tamil for an essay, then the Tamil for a title.
\end{quote}

\subsection*{You'll want to know: \ta{முன்னுரை}}
\textbf{\ta{முன்னுரை}} — \emph{muṉṉurai} — \textquotedblleft{}a preface; an introduction\textquotedblright{}.

\begin{quote}
\ta{இந்தப்} \ta{புத்தகத்தின்} \ta{முன்னுரை} \ta{மிக} \ta{நீளமாக} \ta{இருக்கிறது}.
\end{quote}

\begin{quote}
\emph{intap puttakattiṉ muṉṉurai mika nīḷamāka irukkiṟatu.}
\end{quote}

\begin{quote}
The preface of this book is very long.
\end{quote}

Its partner at the end of an essay is \ta{முடிவுரை}, a conclusion.

\begin{cousinweb}
\textbf{\ta{முன்னுரை}} (\emph{muṉṉurai}) is native: \emph{muṉ}, \textquotedblleft{}in front, before\textquotedblright{} (DEDR 5020), and \emph{urai}, \textquotedblleft{}word, speech\textquotedblright{} (DEDR 648): the words that come first.
\end{cousinweb}

\begin{grammarlens}[title={What you've built}]
One more word for giving reasons and opinions.
\end{grammarlens}

\subsection*{Your turn}
\begin{itemize}
\raggedright
  \item \emph{Say it:} \emph{muṉṉurai}
  \item \emph{Say it:} \emph{muṉṉurai}, once more
  \item \emph{Say it:} \emph{intap puttakattiṉ muṉṉurai mika nīḷamāka irukkiṟatu.} — The preface of this book is very long.
  \item \emph{Recall:} say the Tamil for an essay, then the Tamil for a title, then say \emph{muṉṉurai} again
\end{itemize}

\begin{tcolorbox}[breakable,colback=teal!4,colframe=teal!35!black,title={Before you move on}]
What does \ta{முன்னுரை} mean? (\textquotedblleft{}A preface; an introduction\textquotedblright{}.) Say it once more.
\end{tcolorbox}
\section[\texorpdfstring{\ta{சுருக்கம்} (curukkam)}{சுருக்கம் (curukkam)}]{\ta{சுருக்கம்} (curukkam) — a summary}

\label{lesson:TA-C365-curukkam}



\begin{quote}
Before the new one: say the Tamil for a title, then the Tamil for a preface.
\end{quote}

\subsection*{You'll want to know: \ta{சுருக்கம்}}
\textbf{\ta{சுருக்கம்}} — \emph{curukkam} — \textquotedblleft{}a summary\textquotedblright{}.

\begin{quote}
\ta{அவன்} \ta{கதையின்} \ta{சுருக்கத்தை} \ta{இரண்டு} \ta{வாக்கியங்களில்} \ta{சொன்னான்}.
\end{quote}

\begin{quote}
\emph{avaṉ kataiyiṉ curukkattai iraṇṭu vākkiyaṅkaḷil coṉṉāṉ.}
\end{quote}

\begin{quote}
He gave the summary of the story in two sentences.
\end{quote}

\ta{சுருக்கமாக}, briefly, in short.

\begin{cousinweb}
\textbf{\ta{சுருக்கம்}} (\emph{curukkam}) is native Dravidian: DEDR 2687 lists Tamil \emph{curuṅku}, \textquotedblleft{}to shrink, contract\textquotedblright{}, \emph{curukku}, \textquotedblleft{}to abridge\textquotedblright{}, and \emph{curukkam}, \textquotedblleft{}abridgement\textquotedblright{}.
\end{cousinweb}

\begin{grammarlens}[title={What you've built}]
One more word for giving reasons and opinions.
\end{grammarlens}

\subsection*{Your turn}
\begin{itemize}
\raggedright
  \item \emph{Say it:} \emph{curukkam}
  \item \emph{Say it:} \emph{curukkam}, once more
  \item \emph{Say it:} \emph{avaṉ kataiyiṉ curukkattai iraṇṭu vākkiyaṅkaḷil coṉṉāṉ.} — He gave the summary of the story in two sentences.
  \item \emph{Recall:} say the Tamil for a title, then the Tamil for a preface, then say \emph{curukkam} again
\end{itemize}

\begin{tcolorbox}[breakable,colback=teal!4,colframe=teal!35!black,title={Before you move on}]
What does \ta{சுருக்கம்} mean? (\textquotedblleft{}A summary\textquotedblright{}.) Say it once more.
\end{tcolorbox}
\section[\texorpdfstring{\ta{மேற்கோள்} (mēṟkōḷ)}{மேற்கோள் (mēṟkōḷ)}]{\ta{மேற்கோள்} (mēṟkōḷ) — a quotation}

\label{lesson:TA-C365-merkol}



\begin{quote}
Before the new one: say the Tamil for a preface, then the Tamil for a summary.
\end{quote}

\subsection*{You'll want to know: \ta{மேற்கோள்}}
\textbf{\ta{மேற்கோள்}} — \emph{mēṟkōḷ} — \textquotedblleft{}a quotation\textquotedblright{}.

\begin{quote}
\ta{அவர்} \ta{தன்} \ta{பேச்சில்} \ta{ஒரு} \ta{பழமொழியை} \ta{மேற்கோளாகக்} \ta{காட்டினார்}.
\end{quote}

\begin{quote}
\emph{avar taṉ pēccil oru paḻamoḻiyai mēṟkōḷākak kāṭṭiṉār.}
\end{quote}

\begin{quote}
He quoted a proverb in his speech.
\end{quote}

\ta{மேற்கோள்} \ta{காட்டு}, to quote.

\begin{cousinweb}
\textbf{\ta{மேற்கோள்}} (\emph{mēṟkōḷ}) is native: \emph{mēl}, \textquotedblleft{}that which is above, superior\textquotedblright{} (DEDR 5086), and \emph{kōḷ}, \textquotedblleft{}taking, holding\textquotedblright{} (DEDR 2151); the l of \emph{mēl} becomes \ta{ற்} before the k.
\end{cousinweb}

\begin{grammarlens}[title={What you've built}]
One more word for giving reasons and opinions.
\end{grammarlens}

\subsection*{Your turn}
\begin{itemize}
\raggedright
  \item \emph{Say it:} \emph{mēṟkōḷ}
  \item \emph{Say it:} \emph{mēṟkōḷ}, once more
  \item \emph{Say it:} \emph{avar taṉ pēccil oru paḻamoḻiyai mēṟkōḷākak kāṭṭiṉār.} — He quoted a proverb in his speech.
  \item \emph{Recall:} say the Tamil for a preface, then the Tamil for a summary, then say \emph{mēṟkōḷ} again
\end{itemize}

\begin{tcolorbox}[breakable,colback=teal!4,colframe=teal!35!black,title={Before you move on}]
What does \ta{மேற்கோள்} mean? (\textquotedblleft{}A quotation\textquotedblright{}.) Say it once more.
\end{tcolorbox}
